\section{Zero-frequency flatness and the recorded zero-slack condition}
\label{sec:zero-slack}

We now combine an amplitude estimate with determinant-energy survival.
The calculation is elementary, but its direction is decisive.

\begin{theorem}[General exponent transfer]
\label{thm:general-exponent-transfer}
Fix \(h_0\ne0\).  Assume
\begin{equation}
 Q_X=X^{\beta+o(1)},
 \qquad
 J_X=X^{1-\beta+o(1)}
 \label{eq:general-scales}
\end{equation}
for a fixed \(0<\beta<1\).  Suppose that for fixed
\(\eta_Z,\lambda_D\ge0\),
\begin{align}
 |Z_X|
 &\le X^{-\eta_Z+o(1)}Q_X^2,
 \label{eq:zero-mode-exponent}\\
 D_X
 &\ge X^{-\lambda_D+o(1)}Q_X^3J_X^2.
 \label{eq:energy-exponent}
\end{align}
Then
\begin{equation}
 \boxed{\displaystyle
 \mathfrak F_{h_0}(X)
 \le
 X^{3\beta-2+\lambda_D-2\eta_Z+o(1)}.}
 \label{eq:general-flatness-transfer}
\end{equation}
\end{theorem}

\begin{proof}
The lower bound \eqref{eq:energy-exponent} is positive for all
sufficiently large \(X\), so \(D_X>0\) there.  Hence
\begin{align*}
 \mathfrak F_{h_0}(X)
 &=\frac{|Z_X|^2}{D_X}\\
 &\le
 X^{\lambda_D-2\eta_Z+o(1)}
 \frac{Q_X^4}{Q_X^3J_X^2}\\
 &=X^{\lambda_D-2\eta_Z+o(1)}
 \frac{Q_X}{J_X^2}.
\end{align*}
By \eqref{eq:general-scales},
\(Q_X/J_X^2=X^{\beta-2(1-\beta)+o(1)}
=X^{3\beta-2+o(1)}\).
\end{proof}

\begin{corollary}[The high-row zero-slack transfer]
\label{cor:zero-slack}
At
\begin{equation}
 \beta=\frac{267}{400},
 \qquad
 Q_X=X^{267/400+o(1)},
 \qquad
 J_X=X^{133/400+o(1)},
 \label{eq:high-row-scales}
\end{equation}
the conclusion is
\begin{equation}
 \boxed{\displaystyle
 \mathfrak F_{h_0}(X)
 \le
 X^{1/400+\lambda_D-2\eta_Z+o(1)}.}
 \label{eq:high-row-flatness}
\end{equation}
Consequently, the two recorded exponent bounds imply the TPC-32
hypothesis
\(\mathfrak F_{h_0}(X)\le X^{1/400+o(1)}\) whenever
\begin{equation}
 \boxed{2\eta_Z\ge\lambda_D.}
 \label{eq:zero-slack-condition}
\end{equation}
Strict inequality gives a fixed power of robustness; equality gives
no spare fixed-power margin.  The condition is sharp for a deduction
from the recorded bounds alone; it is not asserted to be necessary
for the physical packet unless the two exponents are optimal.
\end{corollary}

\begin{proof}
The identity
\[
 3\cdot\frac{267}{400}-2
 =\frac{801-800}{400}
 =\frac1{400}
\]
and \cref{thm:general-exponent-transfer} give
\eqref{eq:high-row-flatness}.  Its recorded exponent is at most
\(1/400+o(1)\) under \eqref{eq:zero-slack-condition}.  If that
condition fails, these bounds alone do not close the threshold.
\end{proof}

\begin{corollary}[The current amplitude target requires natural energy]
\label{cor:current-target}
The TPC-33 amplitude target
\[
 |Z_X|\le X^{o(1)}Q_X^2
\]
has \(\eta_Z=0\).  Under that target,
\eqref{eq:zero-slack-condition} permits no positive certified energy
loss.  The sufficient zero-loss requirement is
\begin{equation}
 \lambda_D=0,
 \qquad\text{equivalently}\qquad
 D_X\ge X^{-o(1)}Q_X^3J_X^2.
 \label{eq:natural-energy-required}
\end{equation}
Thus any positive fixed-power determinant-energy loss cannot close
this transfer from the current amplitude bound unless a new
zero-mode saving compensates for it.
\end{corollary}

\begin{proof}
For the sufficient ledger condition, both exponents are nonnegative,
so \(0\ge\lambda_D\) implies \(\lambda_D=0\).
\end{proof}

\begin{remark}[Additional losses]
If a later comparison contributes an additional squared-energy loss
\(\lambda_{\rm aux}\ge0\), then the condition for this augmented
certificate becomes
\[
 2\eta_Z\ge\lambda_D+\lambda_{\rm aux}.
\]
Such a loss cannot be hidden in an \(o(1)\) term when it is a fixed
positive power.
\end{remark}

\begin{remark}[Two distinct appearances of \(1/400\)]
The exponent \(1/400\) in \eqref{eq:high-row-flatness} is the
geometric baseline \(Q_X/J_X^2\) and is the allowed TPC-32 flatness
endpoint.  It is not a free account from which unrecorded losses may
be spent.  Separately, whenever the wider route compares accumulated
fixed-power losses with an inherited \(1/400\) margin, the operational
rule is strict:
\[
 \text{continue only if total loss}<\frac1{400};
 \qquad
 \text{stop if total loss}\ge\frac1{400}.
\]
In particular equality in \eqref{eq:zero-slack-condition} is a valid
conditional implication but has zero robustness against any further
fixed-power degradation.
\end{remark}
